% --- Archon physics formalization source begin ---
% archon:physics
% archon:covers PhyXMiniProblems/problem_phyx_mini_0344.lean
% archon:source-report reports/phyx_mini/problem_phyx_mini_0344.source.json

\chapter{Physics problem phyx\_mini\_0344}
\label{ch:PhyXMiniProblems_problem_phyx_mini_0344}

\paragraph{Problem source.}
\#\# Physical scenario

One-third of a mole of He gas is taken along the path \$abc\$ shown in \textbackslash{}textbf\{figure\}. Assume that the gas may be treated as ideal.

\#\# Question

How much heat is transferred into or out of the gas?

\#\# Answer choices

- A: A: 1.65 \textbackslash{}times 10\textasciicircum{}\{3\}J

- B: B: 3 \textbackslash{}times 10\textasciicircum{}\{3\}J

- C: C: 1.95 \textbackslash{}times 10\textasciicircum{}\{3\}J

- D: D: -1.45 \textbackslash{}times 10\textasciicircum{}\{3\}J

\#\# Figure caption (auxiliary; use the image as primary evidence)

The image is a pressure-volume (p-V) diagram with the following elements:

1. **Axes**: 
   - The vertical axis (y-axis) represents pressure (p) in Pascal (Pa), ranging from \textbackslash{}(1 \textbackslash{}times 10\textasciicircum{}5\textbackslash{}) to \textbackslash{}(4 \textbackslash{}times 10\textasciicircum{}5\textbackslash{}).
   - The horizontal axis (x-axis) represents volume (V) in cubic meters (\textbackslash{}(m\textasciicircum{}3\textbackslash{})), ranging from 0.002 to 0.010.

2. **Points and Path**:
   - Three points labeled \textbackslash{}(a\textbackslash{}), \textbackslash{}(b\textbackslash{}), and \textbackslash{}(c\textbackslash{}) mark significant states in the cyclic process.
   - The path \textbackslash{}(a \textbackslash{}rightarrow b\textbackslash{}) is shown with a blue arrow, indicating an increase in both pressure and volume.
   - The path \textbackslash{}(b \textbackslash{}rightarrow c\textbackslash{}) is shown with a purple arrow, indicating a decrease in pressure and an increase in volume.
   - The path \textbackslash{}(c \textbackslash{}rightarrow a\textbackslash{}) is a dashed line at constant pressure, completing a cycle with decreasing volume.

3. **Text**: 
   - The labels on the axes: \textbackslash{}(p (\textbackslash{}text\{Pa\})\textbackslash{}) for pressure and \textbackslash{}(V (\textbackslash{}text\{m\}\textasciicircum{}3)\textbackslash{}) for volume.
   - The points \textbackslash{}(a\textbackslash{}), \textbackslash{}(b\textbackslash{}), and \textbackslash{}(c\textbackslash{}) are labeled and connected sequentially by the arrows.

This diagram is often used to represent a thermodynamic cycle.

\#\# Dataset metadata

Laws of Thermodynamics; Physical Model Grounding Reasoning, Multi-Formula Reasoning

\paragraph{Recorded answer/context.}
B: B: 3 \textbackslash{}times 10\textasciicircum{}\{3\}J

\paragraph{Figure/image path.}
/root/proposal\_for\_physic/science-mango/phyx\_mini\_run/phyx\_data/test\_image/344.png

\paragraph{Formalization target.}
create a compiling Lean file with sorry bodies at `PhyXMiniProblems/problem\_phyx\_mini\_0344.lean`. The Lean declarations must preserve the physical quantities, dimensions or dimensional roles, figure labels, governing-law hypotheses, and final relation expressed by this problem.
Use Mathlib/Physlib names found through LeanExplore where available. If a physics API is missing, introduce faithful local abstractions rather than scalar placeholder aliases.

\begin{theorem}[Physics formalization target]
\label{thm:physics:phyx_mini_0344:target}
\lean{PhyXMiniProblems.ProblemPhyXMini0344.problem_phyx_mini_0344}
\uses{def:physics:phyx-mini-0344:phyxminiproblems-problemphyxmini0344-energyinjoules, def:physics:phyx-mini-0344:phyxminiproblems-problemphyxmini0344-heliumpvprocess, def:physics:phyx-mini-0344:phyxminiproblems-problemphyxmini0344-matchesproblemstatementandfigure, def:physics:phyx-mini-0344:phyxminiproblems-problemphyxmini0344-hasphysicalgasparameters, def:physics:phyx-mini-0344:phyxminiproblems-problemphyxmini0344-satisfiesidealmonatomicgaslaws, def:physics:phyx-mini-0344:phyxminiproblems-problemphyxmini0344-satisfiesstraightpvworklaw, def:physics:phyx-mini-0344:phyxminiproblems-problemphyxmini0344-satisfiesfirstlawalongabc, def:physics:phyx-mini-0344:phyxminiproblems-problemphyxmini0344-answerheatinjoules, def:physics:phyx-mini-0344:phyxminiproblems-problemphyxmini0344-recordedanswerchoice}
The assigned autoformalize agent should translate this physics problem into Lean declarations in the covered file, with theorem and lemma proof bodies written as `by sorry`.
\end{theorem}
\begin{proof}
This is an autoformalization task, not a proof task. Produce statements that can later be proved by the physics prover without weakening the physical model.
\end{proof}

% --- Archon declaration topology begin ---
\section{Declaration topology}
\begin{definition}[Volume Quantity]
  \label{def:physics:phyx-mini-0344:phyxminiproblems-problemphyxmini0344-volumequantity}
  \lean{PhyXMiniProblems.ProblemPhyXMini0344.VolumeQuantity}
  A nonnegative physical volume, with dimension length cubed
\end{definition}
\begin{proof}
  This is the declared model definition of Volume Quantity.
\end{proof}
\begin{definition}[Pressure Quantity]
  \label{def:physics:phyx-mini-0344:phyxminiproblems-problemphyxmini0344-pressurequantity}
  \lean{PhyXMiniProblems.ProblemPhyXMini0344.PressureQuantity}
  A physical pressure, represented by Physlib's pressure quantity
\end{definition}
\begin{proof}
  This is the declared model definition of Pressure Quantity.
\end{proof}
\begin{definition}[Energy Quantity]
  \label{def:physics:phyx-mini-0344:phyxminiproblems-problemphyxmini0344-energyquantity}
  \lean{PhyXMiniProblems.ProblemPhyXMini0344.EnergyQuantity}
  A physical energy; its sign records the direction of an energy transfer
\end{definition}
\begin{proof}
  This is the declared model definition of Energy Quantity.
\end{proof}
\begin{definition}[volume In Cubic Meters]
  \label{def:physics:phyx-mini-0344:phyxminiproblems-problemphyxmini0344-volumeincubicmeters}
  \lean{PhyXMiniProblems.ProblemPhyXMini0344.volumeInCubicMeters}
  \uses{def:physics:phyx-mini-0344:phyxminiproblems-problemphyxmini0344-volumequantity}
  SI volume readout, in cubic metres
\end{definition}
\begin{proof}
  This is the declared model definition of volume In Cubic Meters.
\end{proof}
\begin{definition}[pressure In Pascals]
  \label{def:physics:phyx-mini-0344:phyxminiproblems-problemphyxmini0344-pressureinpascals}
  \lean{PhyXMiniProblems.ProblemPhyXMini0344.pressureInPascals}
  \uses{def:physics:phyx-mini-0344:phyxminiproblems-problemphyxmini0344-pressurequantity}
  SI pressure readout, in pascals
\end{definition}
\begin{proof}
  This is the declared model definition of pressure In Pascals.
\end{proof}
\begin{definition}[energy In Joules]
  \label{def:physics:phyx-mini-0344:phyxminiproblems-problemphyxmini0344-energyinjoules}
  \lean{PhyXMiniProblems.ProblemPhyXMini0344.energyInJoules}
  \uses{def:physics:phyx-mini-0344:phyxminiproblems-problemphyxmini0344-energyquantity}
  SI energy readout, in joules
\end{definition}
\begin{proof}
  This is the declared model definition of energy In Joules.
\end{proof}
\begin{definition}[Gas State]
  \label{def:physics:phyx-mini-0344:phyxminiproblems-problemphyxmini0344-gasstate}
  \lean{PhyXMiniProblems.ProblemPhyXMini0344.GasState}
  The three labeled equilibrium states in the supplied p-V diagram
\end{definition}
\begin{proof}
  This is the declared model definition of Gas State.
\end{proof}
\begin{definition}[Process Leg]
  \label{def:physics:phyx-mini-0344:phyxminiproblems-problemphyxmini0344-processleg}
  \lean{PhyXMiniProblems.ProblemPhyXMini0344.ProcessLeg}
  The three drawn connections between labeled states
\end{definition}
\begin{proof}
  This is the declared model definition of Process Leg.
\end{proof}
\begin{definition}[Gas Species]
  \label{def:physics:phyx-mini-0344:phyxminiproblems-problemphyxmini0344-gasspecies}
  \lean{PhyXMiniProblems.ProblemPhyXMini0344.GasSpecies}
  The chemical species named in the problem
\end{definition}
\begin{proof}
  This is the declared model definition of Gas Species.
\end{proof}
\begin{definition}[Gas Model]
  \label{def:physics:phyx-mini-0344:phyxminiproblems-problemphyxmini0344-gasmodel}
  \lean{PhyXMiniProblems.ProblemPhyXMini0344.GasModel}
  The thermodynamic model requested in the problem statement
\end{definition}
\begin{proof}
  This is the declared model definition of Gas Model.
\end{proof}
\begin{definition}[Process Regime]
  \label{def:physics:phyx-mini-0344:phyxminiproblems-problemphyxmini0344-processregime}
  \lean{PhyXMiniProblems.ProblemPhyXMini0344.ProcessRegime}
  The equilibrium-process regime needed to interpret the drawn p-V path
\end{definition}
\begin{proof}
  This is the declared model definition of Process Regime.
\end{proof}
\begin{definition}[Figure Axis]
  \label{def:physics:phyx-mini-0344:phyxminiproblems-problemphyxmini0344-figureaxis}
  \lean{PhyXMiniProblems.ProblemPhyXMini0344.FigureAxis}
  The two axes of the primary diagram
\end{definition}
\begin{proof}
  This is the declared model definition of Figure Axis.
\end{proof}
\begin{definition}[Axis Role]
  \label{def:physics:phyx-mini-0344:phyxminiproblems-problemphyxmini0344-axisrole}
  \lean{PhyXMiniProblems.ProblemPhyXMini0344.AxisRole}
  Physical roles of the plotted axes
\end{definition}
\begin{proof}
  This is the declared model definition of Axis Role.
\end{proof}
\begin{definition}[Axis Display Unit]
  \label{def:physics:phyx-mini-0344:phyxminiproblems-problemphyxmini0344-axisdisplayunit}
  \lean{PhyXMiniProblems.ProblemPhyXMini0344.AxisDisplayUnit}
  Units printed beside the axes
\end{definition}
\begin{proof}
  This is the declared model definition of Axis Display Unit.
\end{proof}
\begin{definition}[Figure Label]
  \label{def:physics:phyx-mini-0344:phyxminiproblems-problemphyxmini0344-figurelabel}
  \lean{PhyXMiniProblems.ProblemPhyXMini0344.FigureLabel}
  Mathematical labels printed in the diagram
\end{definition}
\begin{proof}
  This is the declared model definition of Figure Label.
\end{proof}
\begin{definition}[Leg Style]
  \label{def:physics:phyx-mini-0344:phyxminiproblems-problemphyxmini0344-legstyle}
  \lean{PhyXMiniProblems.ProblemPhyXMini0344.LegStyle}
  Appearance of a process leg in the supplied bitmap
\end{definition}
\begin{proof}
  This is the declared model definition of Leg Style.
\end{proof}
\begin{definition}[Leg Geometry]
  \label{def:physics:phyx-mini-0344:phyxminiproblems-problemphyxmini0344-leggeometry}
  \lean{PhyXMiniProblems.ProblemPhyXMini0344.LegGeometry}
  Geometric shape of a plotted process leg
\end{definition}
\begin{proof}
  This is the declared model definition of Leg Geometry.
\end{proof}
\begin{definition}[Pressure Volume Figure]
  \label{def:physics:phyx-mini-0344:phyxminiproblems-problemphyxmini0344-pressurevolumefigure}
  \lean{PhyXMiniProblems.ProblemPhyXMini0344.PressureVolumeFigure}
  \uses{def:physics:phyx-mini-0344:phyxminiproblems-problemphyxmini0344-volumequantity, def:physics:phyx-mini-0344:phyxminiproblems-problemphyxmini0344-pressurequantity, def:physics:phyx-mini-0344:phyxminiproblems-problemphyxmini0344-gasstate, def:physics:phyx-mini-0344:phyxminiproblems-problemphyxmini0344-processleg, def:physics:phyx-mini-0344:phyxminiproblems-problemphyxmini0344-figureaxis, def:physics:phyx-mini-0344:phyxminiproblems-problemphyxmini0344-axisrole, def:physics:phyx-mini-0344:phyxminiproblems-problemphyxmini0344-axisdisplayunit, def:physics:phyx-mini-0344:phyxminiproblems-problemphyxmini0344-figurelabel, def:physics:phyx-mini-0344:phyxminiproblems-problemphyxmini0344-legstyle, def:physics:phyx-mini-0344:phyxminiproblems-problemphyxmini0344-leggeometry}
  The diagram as an independent source of state coordinates and path metadata. Keeping these coordinates separate from the gas state lets the data predicate say explicitly that the physical process passes through the plotted points
\end{definition}
\begin{proof}
  This is the declared model definition of Pressure Volume Figure.
\end{proof}
\begin{definition}[Helium PVProcess]
  \label{def:physics:phyx-mini-0344:phyxminiproblems-problemphyxmini0344-heliumpvprocess}
  \lean{PhyXMiniProblems.ProblemPhyXMini0344.HeliumPVProcess}
  \uses{def:physics:phyx-mini-0344:phyxminiproblems-problemphyxmini0344-volumequantity, def:physics:phyx-mini-0344:phyxminiproblems-problemphyxmini0344-pressurequantity, def:physics:phyx-mini-0344:phyxminiproblems-problemphyxmini0344-energyquantity, def:physics:phyx-mini-0344:phyxminiproblems-problemphyxmini0344-gasstate, def:physics:phyx-mini-0344:phyxminiproblems-problemphyxmini0344-processleg, def:physics:phyx-mini-0344:phyxminiproblems-problemphyxmini0344-gasspecies, def:physics:phyx-mini-0344:phyxminiproblems-problemphyxmini0344-gasmodel, def:physics:phyx-mini-0344:phyxminiproblems-problemphyxmini0344-processregime, def:physics:phyx-mini-0344:phyxminiproblems-problemphyxmini0344-pressurevolumefigure}
  The physical helium sample and its path-dependent energy transfers. amountOfGasMoles and molarGasConstantJoulesPerMoleKelvin are named scalar readouts because Physlib's unit system has no amount-of-substance dimension. The signed heat convention is positive for heat transferred into the gas, and work is positive when done by the gas on its environment. No field is defined from the requested numerical heat
\end{definition}
\begin{proof}
  This is the declared model definition of Helium PVProcess.
\end{proof}
\begin{definition}[Matches Problem Statement And Figure]
  \label{def:physics:phyx-mini-0344:phyxminiproblems-problemphyxmini0344-matchesproblemstatementandfigure}
  \lean{PhyXMiniProblems.ProblemPhyXMini0344.MatchesProblemStatementAndFigure}
  \uses{def:physics:phyx-mini-0344:phyxminiproblems-problemphyxmini0344-volumeincubicmeters, def:physics:phyx-mini-0344:phyxminiproblems-problemphyxmini0344-pressureinpascals, def:physics:phyx-mini-0344:phyxminiproblems-problemphyxmini0344-gasstate, def:physics:phyx-mini-0344:phyxminiproblems-problemphyxmini0344-processleg, def:physics:phyx-mini-0344:phyxminiproblems-problemphyxmini0344-heliumpvprocess}
  Facts stated in the prose or transcribed from the primary bitmap. The point b lies halfway between the 3 × 10⁵ Pa and 4 × 10⁵ Pa ticks. The two solid directed legs form the questioned path; the dashed c-to-a leg merely shows the closure. None of these fields mentions heat, work, or internal energy
\end{definition}
\begin{proof}
  This is the declared model definition of Matches Problem Statement And Figure.
\end{proof}
\begin{definition}[Has Physical Gas Parameters]
  \label{def:physics:phyx-mini-0344:phyxminiproblems-problemphyxmini0344-hasphysicalgasparameters}
  \lean{PhyXMiniProblems.ProblemPhyXMini0344.HasPhysicalGasParameters}
  \uses{def:physics:phyx-mini-0344:phyxminiproblems-problemphyxmini0344-volumeincubicmeters, def:physics:phyx-mini-0344:phyxminiproblems-problemphyxmini0344-pressureinpascals, def:physics:phyx-mini-0344:phyxminiproblems-problemphyxmini0344-gasstate, def:physics:phyx-mini-0344:phyxminiproblems-problemphyxmini0344-heliumpvprocess}
  Positivity and nondegeneracy assumptions selecting a physical ideal-gas process. In particular, no sign condition is imposed on the requested heat or on the derived work
\end{definition}
\begin{proof}
  This is the declared model definition of Has Physical Gas Parameters.
\end{proof}
\begin{definition}[Satisfies Ideal Monatomic Gas Laws]
  \label{def:physics:phyx-mini-0344:phyxminiproblems-problemphyxmini0344-satisfiesidealmonatomicgaslaws}
  \lean{PhyXMiniProblems.ProblemPhyXMini0344.SatisfiesIdealMonatomicGasLaws}
  \uses{def:physics:phyx-mini-0344:phyxminiproblems-problemphyxmini0344-volumeincubicmeters, def:physics:phyx-mini-0344:phyxminiproblems-problemphyxmini0344-pressureinpascals, def:physics:phyx-mini-0344:phyxminiproblems-problemphyxmini0344-energyinjoules, def:physics:phyx-mini-0344:phyxminiproblems-problemphyxmini0344-gasstate, def:physics:phyx-mini-0344:phyxminiproblems-problemphyxmini0344-heliumpvprocess}
  Macroscopic ideal-gas and monatomic-helium energy laws in coherent SI readouts. The first relation is pV = nRT; the second is U = (3/2)nRT. They apply at every state and contain no heat-transfer conclusion
\end{definition}
\begin{proof}
  This is the declared model definition of Satisfies Ideal Monatomic Gas Laws.
\end{proof}
\begin{definition}[Satisfies Straight PVWork Law]
  \label{def:physics:phyx-mini-0344:phyxminiproblems-problemphyxmini0344-satisfiesstraightpvworklaw}
  \lean{PhyXMiniProblems.ProblemPhyXMini0344.SatisfiesStraightPVWorkLaw}
  \uses{def:physics:phyx-mini-0344:phyxminiproblems-problemphyxmini0344-volumeincubicmeters, def:physics:phyx-mini-0344:phyxminiproblems-problemphyxmini0344-pressureinpascals, def:physics:phyx-mini-0344:phyxminiproblems-problemphyxmini0344-energyinjoules, def:physics:phyx-mini-0344:phyxminiproblems-problemphyxmini0344-processleg, def:physics:phyx-mini-0344:phyxminiproblems-problemphyxmini0344-heliumpvprocess}
  For each straight quasistatic leg on abc, integrating p dV gives the trapezoid formula: average endpoint pressure times the volume change. Work along the questioned path is additive over its two legs. This is a generic work law and does not assume the derived 1800 J value
\end{definition}
\begin{proof}
  This is the declared model definition of Satisfies Straight PVWork Law.
\end{proof}
\begin{definition}[Satisfies First Law Along ABC]
  \label{def:physics:phyx-mini-0344:phyxminiproblems-problemphyxmini0344-satisfiesfirstlawalongabc}
  \lean{PhyXMiniProblems.ProblemPhyXMini0344.SatisfiesFirstLawAlongABC}
  \uses{def:physics:phyx-mini-0344:phyxminiproblems-problemphyxmini0344-energyinjoules, def:physics:phyx-mini-0344:phyxminiproblems-problemphyxmini0344-heliumpvprocess}
  First-law sign convention on the path a → b → c: Q\_into = (U\_c - U\_a) + W\_by\_gas. This governing relation leaves every term as an independent physical energy; it does not insert the requested numerical value or its sign
\end{definition}
\begin{proof}
  This is the declared model definition of Satisfies First Law Along ABC.
\end{proof}
\begin{lemma}[internal Energy Change Along ABC eq 1200 joules]
  \label{lem:physics:phyx-mini-0344:phyxminiproblems-problemphyxmini0344-internalenergychangealongabc-eq-1200-joules}
  \lean{PhyXMiniProblems.ProblemPhyXMini0344.internalEnergyChangeAlongABC_eq_1200_joules}
  \uses{def:physics:phyx-mini-0344:phyxminiproblems-problemphyxmini0344-energyinjoules, def:physics:phyx-mini-0344:phyxminiproblems-problemphyxmini0344-heliumpvprocess, def:physics:phyx-mini-0344:phyxminiproblems-problemphyxmini0344-matchesproblemstatementandfigure, def:physics:phyx-mini-0344:phyxminiproblems-problemphyxmini0344-hasphysicalgasparameters, def:physics:phyx-mini-0344:phyxminiproblems-problemphyxmini0344-satisfiesidealmonatomicgaslaws}
  The endpoint products and the monatomic law give ΔU = 1200 J
\end{lemma}
\begin{proof}
  The conclusion follows from the governing relations and figure readouts named in the statement of internal Energy Change Along ABC eq 1200 joules.
\end{proof}
\begin{lemma}[work Done By Gas Along ABC eq 1800 joules]
  \label{lem:physics:phyx-mini-0344:phyxminiproblems-problemphyxmini0344-workdonebygasalongabc-eq-1800-joules}
  \lean{PhyXMiniProblems.ProblemPhyXMini0344.workDoneByGasAlongABC_eq_1800_joules}
  \uses{def:physics:phyx-mini-0344:phyxminiproblems-problemphyxmini0344-energyinjoules, def:physics:phyx-mini-0344:phyxminiproblems-problemphyxmini0344-heliumpvprocess, def:physics:phyx-mini-0344:phyxminiproblems-problemphyxmini0344-matchesproblemstatementandfigure, def:physics:phyx-mini-0344:phyxminiproblems-problemphyxmini0344-satisfiesstraightpvworklaw}
  The two straight-line trapezoids give W\_by\_gas = 1800 J
\end{lemma}
\begin{proof}
  The conclusion follows from the governing relations and figure readouts named in the statement of work Done By Gas Along ABC eq 1800 joules.
\end{proof}
\begin{definition}[Answer Choice]
  \label{def:physics:phyx-mini-0344:phyxminiproblems-problemphyxmini0344-answerchoice}
  \lean{PhyXMiniProblems.ProblemPhyXMini0344.AnswerChoice}
  Labels of the four displayed heat-transfer choices
\end{definition}
\begin{proof}
  This is the declared model definition of Answer Choice.
\end{proof}
\begin{definition}[answer Heat In Joules]
  \label{def:physics:phyx-mini-0344:phyxminiproblems-problemphyxmini0344-answerheatinjoules}
  \lean{PhyXMiniProblems.ProblemPhyXMini0344.answerHeatInJoules}
  \uses{def:physics:phyx-mini-0344:phyxminiproblems-problemphyxmini0344-answerchoice}
  Signed heat readout, in joules, printed beside each answer label
\end{definition}
\begin{proof}
  This is the declared model definition of answer Heat In Joules.
\end{proof}
\begin{definition}[recorded Answer Choice]
  \label{def:physics:phyx-mini-0344:phyxminiproblems-problemphyxmini0344-recordedanswerchoice}
  \lean{PhyXMiniProblems.ProblemPhyXMini0344.recordedAnswerChoice}
  \uses{def:physics:phyx-mini-0344:phyxminiproblems-problemphyxmini0344-answerchoice}
  The answer label recorded in the supplied dataset metadata
\end{definition}
\begin{proof}
  This is the declared model definition of recorded Answer Choice.
\end{proof}
% --- Archon declaration topology end ---
% --- Archon physics formalization source end ---
